\section{Tổng quan tài liệu}

Phần Tổng quan tài liệu sẽ tổng hợp và phân tích các nghiên cứu trước đây liên quan đến ấn phẩm khoa học bị thu hồi, đặc biệt trong lĩnh vực Khoa học thông tin, nhằm xác định các hướng tiếp cận, nguồn dữ liệu, phương pháp phân tích và các khía cạnh đã được quan tâm trong nghiên cứu hiện có. Nội dung tổng quan có thể tập trung vào các nghiên cứu sử dụng trắc lượng thư mục và phân tích dữ liệu để khảo sát xu hướng thu hồi, đặc điểm của tác giả, quốc gia, tổ chức, tạp chí, chủ đề nghiên cứu, nguyên nhân thu hồi và mạng lưới hợp tác khoa học. Bên cạnh đó, phần Tổng quan tài liệu sẽ so sánh các nguồn dữ liệu thư mục và cơ sở dữ liệu về việc thu hồi ấn phẩm, qua đó làm rõ ưu điểm, hạn chế và khả năng kết hợp dữ liệu từ nhiều nguồn. Từ đó, sẽ chỉ ra khoảng trống nghiên cứu và làm rõ cách nghiên cứu hiện tại bổ sung cho các công trình trước đây.